\documentclass[paper=a4, fontsize=11pt]{jlreq} % LuaLaTeX または upLaTeX でのコンパイルを推奨

% --- パッケージ群 ---
\usepackage[margin=20mm]{geometry} % 上下左右の余白を20mmに設定
\usepackage{amsmath, amssymb, amsthm, amsfonts} % 数学の基本
\usepackage{physics} % \ket{0}, \pdv{x} などが簡単に打てる
\usepackage{bm} % 太字ベクトル \bm{v}
\usepackage{graphicx}
\usepackage{hyperref} % リンク・目次用
\usepackage{cite} % 引用用
\usepackage{etoolbox} % 環境の終わりに自動でマークを入れるために追加
\usepackage{ytableau} % ★追加: Young図形の描画に必須
\usepackage{tikz}     % ★追加: TikZの読み込み
\usepackage{mathtools}
\usepackage[notref,notcite]{showkeys}
\usepackage{xcolor}
\usepackage{mdframed}
\usetikzlibrary{arrows.meta, decorations.pathreplacing}% ★追加: TikZの波括弧や矢印のスタイルに必須

% 本文用の環境（色は目に優しいDarkBlueなどを指定）
\newenvironment{story}{%
  \begin{mdframed}[
    backgroundcolor=red!5, % redを5%の濃さで出力（ほんのり薄い赤）
    linewidth=0pt,         % 枠線はなし（スッキリさせるため）
    innertopmargin=10pt,   % 上下の余白
    innerbottommargin=10pt,
    innerrightmargin=10pt, % 左右の余白
    innerleftmargin=10pt
  ]%
}{%
  \end{mdframed}%
}

\newtheoremstyle{break}% スタイル名
  {1em}% 上部のスペース（約1行分空ける）
  {1em}% 下部のスペース（約1行分空ける）
  {\normalfont}% 本文のフォント（日本語なので斜体にしない）
  {}% インデント量
  {\bfseries}% 見出しのフォント（太字）
  {}% 見出し後の句読点（不要なら空に）
  {\newline}% ★ここで改行を指定！★
  {}% 見出しのカスタム指定（デフォルト）

% --- 定理環境の設定 ---
\theoremstyle{break}
\newtheorem{theorem}{定理}[section]
\newtheorem{definition}[theorem]{定義}
\newtheorem{proposition}[theorem]{命題}
\newtheorem{lemma}[theorem]{補題}
\newtheorem{corollary}[theorem]{系}
\newtheorem{example}[theorem]{例}
\newtheorem{problem}[theorem]{問題} % 自作問題用


% --- 環境の終わりにマークをつける設定 ---
\newcommand{\myendmark}{\null\hfill$\diamondsuit$} % 既存のコマンドと被らないよう \myendmark に変更
\AtEndEnvironment{theorem}{\myendmark}
\AtEndEnvironment{definition}{\myendmark}
\AtEndEnvironment{proposition}{\myendmark}
\AtEndEnvironment{lemma}{\myendmark}
\AtEndEnvironment{corollary}{\myendmark}
\AtEndEnvironment{example}{\myendmark}
\AtEndEnvironment{problem}{\myendmark}

% --- Q&A（自問自答）用の専用スタイル ---
\newtheoremstyle{qa}% スタイル名
  {1ex}% 上部のスペース
  {1ex}% 下部のスペース
  {\normalfont}% 本文のフォント（日本語が崩れないようnormalfontを指定）
  {}% インデント量
  {\bfseries}% 見出しのフォント（太字）
  {}% 見出し後の句読点
  {\newline}% 見出しの後に改行を入れる
  {\thmname{#1}\thmnote{：#3}}% カスタム指定（例：「Q：なぜ〜？」と出力）

\theoremstyle{qa}
\newtheorem*{Q}{Q}

\begin{document}

\begin{story}
以下、前節から考察している2階フックス型線型常微分方程式
\begin{equation}
    \frac{d^2 y}{dx^2} + p_1(x,t) \frac{dy}{dx} + p_2(x,t)y = 0 \label{eq:3.33}
\end{equation}
のモノドロミー保存変形の計算に戻ろう．係数は
\end{story}
\begin{align}
p_1(x,t) &= \frac{1-\kappa_0}{x} + \frac{1-\kappa_1}{x-1} + \sum_{l=1}^g \frac{1-\theta_l}{x-t_l} - \sum_{k=1}^g \frac{1}{x-\lambda_k} \label{eq:3.67} \\
p_2(x,t) &= \frac{\kappa}{x(x-1)} + \sum_{k=1}^g \frac{\lambda_k(\lambda_k-1)\mu_k}{x(x-1)(x-\lambda_k)} - \sum_{l=1}^g \frac{t_l(t_l-1)H_l}{x(x-1)(x-t_l)} \label{eq:3.68}
\end{align}
\begin{story}
により与えられていた．我々は，$x=\lambda_k$ が見かけの特異点である，という仮定を使って，$H_l$ を他の変数の有理関数として具体的に求めた。
\end{story}
\begin{align}
H_l &= M_l \left[ \sum_{k=1}^g M^{k,l} \left\{ \mu_k^2 - A_{k,l}\mu_k + \frac{\kappa}{\lambda_k(\lambda_k-1)} \right\} \right] \label{eq:3.73}\\
A_{k,l} &= \frac{\kappa_0}{\lambda_k} + \frac{\kappa_1}{\lambda_k-1} + \sum_{h=1}^g \frac{\theta_h - \delta_{hl}}{\lambda_k - t_h}\\
M_l &= - \frac{\mathrm{L}(t_l)}{\mathrm{T}'(t_l)} \label{eq:3.74} \\
M^{k,l} &= \frac{\mathrm{T}(\lambda_k)}{\mathrm{L}'(\lambda_k)(\lambda_k - t_l)} \label{eq:3.75}
\end{align}
\begin{story}
ここで、第3.2節において調べたことを思いだそう．(3.33)に対して
\begin{equation}
    r(x,t) = -p_2(x,t) + \frac{1}{4}p_1(x,t)^2 + \frac{1}{2}\frac{\partial}{\partial x}p_1(x,t) \label{eq:3.32}
\end{equation}
とし、SL型方程式
\begin{equation}
    \frac{d^2 z}{dx^2} = r(x,t)z
\end{equation}
を併せて考える．ここで124ページの考察を補題の形にまとめておこう．
\end{story}
124ページの命題3.7では
\[
\frac{d\phi}{dx}+\frac{1}{2}p_{1}(x,t)\phi=0
\]
のモノドロミーが$t$によらないことを仮定していたが、すぐ後で$\mathbb{P}^{1}(\mathbb{C})$では自動的にその条件が満たされることを言っている。具体的に解くと、
\[
\phi(x,t)=\prod_{i=1}^{M}(x-\xi_{i})^{-\frac{1}{2}\alpha_{i}}
\]
で与えられてモノドロミーは$e^{2\pi i\alpha_{i}(t)}$というスカラーで書けている。
\[
\frac{\partial^{2}\vec{\phi}}{\partial x^{2}} + p_{1}(x,t)\frac{\partial\vec{\phi}}{\partial x}+p_{2}(x,t)\vec{\phi}
\]
のモノドロミーが$t$によらないならモノドロミー行列は対角行列で、その成分は、$x=\xi$での特性指数を$s_{1},s_{2}$とすると、$e^{2\pi is_{i}}$の形であるから$s_{i}$が$t$によらないことは必要。つまり、$\alpha_{i}$は特性指数の和で書けているので$t$によらず、$\phi$の$x=\xi_{i}$におけるモノドロミー行列は$t$によらない。よって以下が導ける。
\begin{story}
\begin{lemma}[3.5]
\begin{equation}
    \frac{d^2 y}{dx^2} + p_1(x,t)\frac{dy}{dx} + p_2(x,t)y = 0 \label{eq:3.33}
\end{equation}
がモノドロミー保存変形を定めることと，
\begin{equation}
    \frac{d^2 z}{dx^2} = r(x,t)z \label{eq:3.25}
\end{equation}
がそうなることは同値であり，さらにこれは，以下の偏微分方程式系が $x$ の有理関数を解としてもつことと同値である．
\begin{align}
    \frac{\partial^3}{\partial x^3}w_l - 4r(x,t)\frac{\partial}{\partial x}w_l - 2\frac{\partial r}{\partial x}(x,t)w_l + 2\frac{\partial r}{\partial t_l}(x,t) = 0 \label{eq:3.82} \\
    \left( \frac{\partial}{\partial t_h} - w_h \frac{\partial}{\partial x} \right) w_l = \left( \frac{\partial}{\partial t_l} - w_l \frac{\partial}{\partial x} \right) w_h \quad (h,l = 1,\cdots,g) \label{eq:3.83}
\end{align}
\end{lemma}
実際、
\begin{proposition}[3.8]
非斉次微分方程式
\begin{equation}
    \frac{\partial^3 w}{\partial x^3} - 4r\frac{\partial w}{\partial x} - 2\frac{\partial r}{\partial x}w + 2\frac{\partial r}{\partial t_l} = 0 \label{eq:3.29}
\end{equation}
を満たす，$x$ の有理関数は存在するとしてもただ1つである。
\end{proposition}
により(3.82), (3.83)の有理関数解 $w_l = a_2^{(l)}(x,t)$ は存在すれば唯一であり，それに対して $a_1^{(l)}(x,t)$ を
\begin{equation}
    \frac{\partial p_1}{\partial t_l} = \frac{\partial c^{(l)}}{\partial x}, \quad c^{(l)} = p_1 a_2^{(l)} - 2a_1^{(l)} - \frac{\partial a_2^{(l)}}{\partial x} \label{eq:3.23}
\end{equation}
により定めると，変形方程式系
\begin{equation}
\left\{
\begin{array}{l}
    \displaystyle \frac{\partial^2 \vec{\phi}}{\partial x^2} + p_1(x,t)\frac{\partial \vec{\phi}}{\partial x} + p_2(x,t)\vec{\phi} = 0 \vspace{0.5em} \\
    \displaystyle \frac{\partial \vec{\phi}}{\partial t_l} = a_1^{(l)}(x,t)\vec{\phi} + a_2^{(l)}(x,t)\frac{\partial \vec{\phi}}{\partial x}
\end{array}
\right. \label{eq:3.20}
\end{equation}
は完全積分可能である．さらに次の命題も成り立つ．

\begin{proposition}[3.14]
偏微分方程式系(3.82)の有理関数解 $w_l$ は，(3.83)を満たす．
\end{proposition}

この命題により，モノドロミー保存変形は(3.82)の考察に帰着される．証明のために
\begin{equation*}
    L = \frac{\partial^3}{\partial x^3} - 4r(x,t)\frac{\partial}{\partial x} - 2\frac{\partial r}{\partial x}(x,t)
\end{equation*}
とおく．計算によって次の補題の成立を確認することができる．

\begin{lemma}[3.6]
任意の解析関数 $f,g$ に対して次式が成り立つ．
\begin{equation*}
    2\frac{\partial f}{\partial x}L(g) - 2\frac{\partial g}{\partial x}L(f) + f\frac{\partial}{\partial x}L(g) - g\frac{\partial}{\partial x}L(f) = L\left( f\frac{\partial g}{\partial x} - g\frac{\partial f}{\partial x} \right)
\end{equation*}
\end{lemma}
\end{story}
\begin{proof}
\begin{align*}
    & 2 \frac{\partial f}{\partial x} L(g) - 2 \frac{\partial g}{\partial x} L(f) + f \frac{\partial}{\partial x} L(g) - g \frac{\partial}{\partial x} L(f) \\
    &= 2 \frac{\partial f}{\partial x} \left( \frac{\partial^3 g}{\partial x^3} - 4r \frac{\partial g}{\partial x} - 2 \frac{\partial r}{\partial x} g \right) - 2 \frac{\partial g}{\partial x} \left( \frac{\partial^3 f}{\partial x^3} - 4r \frac{\partial f}{\partial x} - 2 \frac{\partial r}{\partial x} f \right) \\
    &\quad + f \frac{\partial}{\partial x} \left( \frac{\partial^3 g}{\partial x^3} - 4r \frac{\partial g}{\partial x} - 2 \frac{\partial r}{\partial x} g \right) - g \frac{\partial}{\partial x} \left( \frac{\partial^3 f}{\partial x^3} - 4r \frac{\partial f}{\partial x} - 2 \frac{\partial r}{\partial x} f \right) \\
    &= 2 \left( \frac{\partial f}{\partial x}\frac{\partial^3 g}{\partial x^3} - \frac{\partial g}{\partial x}\frac{\partial^3 f}{\partial x^3} \right) - 4 \frac{\partial r}{\partial x} \left( \frac{\partial f}{\partial x}g - \frac{\partial g}{\partial x}f \right) \\
    &\quad + f \left( \frac{\partial^4 g}{\partial x^4} - 4 \frac{\partial r}{\partial x}\frac{\partial g}{\partial x} - 4r \frac{\partial^2 g}{\partial x^2} - 2 \frac{\partial^2 r}{\partial x^2}g - 2 \frac{\partial r}{\partial x}\frac{\partial g}{\partial x} \right) \\
    &\quad - g \left( \frac{\partial^4 f}{\partial x^4} - 4 \frac{\partial r}{\partial x}\frac{\partial f}{\partial x} - 4r \frac{\partial^2 f}{\partial x^2} - 2 \frac{\partial^2 r}{\partial x^2}f - 2 \frac{\partial r}{\partial x}\frac{\partial f}{\partial x} \right) \\
    &= \left( 2\frac{\partial f}{\partial x}\frac{\partial^3 g}{\partial x^3} + f\frac{\partial^4 g}{\partial x^4} - 2\frac{\partial g}{\partial x}\frac{\partial^3 f}{\partial x^3} - g\frac{\partial^4 f}{\partial x^4} \right) - 4r \left( f\frac{\partial^2 g}{\partial x^2} - g\frac{\partial^2 f}{\partial x^2} \right) - 2\frac{\partial r}{\partial x} \left( f\frac{\partial g}{\partial x} - g\frac{\partial f}{\partial x} \right)
\end{align*}

ここで、
\begin{align*}
    \frac{\partial^3}{\partial x^3} \left( f\frac{\partial g}{\partial x} - g\frac{\partial f}{\partial x} \right) &= \frac{\partial^2}{\partial x^2} \left( f\frac{\partial^2 g}{\partial x^2} - g\frac{\partial^2 f}{\partial x^2} \right) \\
    &= \frac{\partial}{\partial x} \left( \frac{\partial f}{\partial x}\frac{\partial^2 g}{\partial x^2} + f\frac{\partial^3 g}{\partial x^3} - \frac{\partial g}{\partial x}\frac{\partial^2 f}{\partial x^2} - g\frac{\partial^3 f}{\partial x^3} \right) \\
    &= \left( 2\frac{\partial f}{\partial x}\frac{\partial^3 g}{\partial x^3} + f\frac{\partial^4 g}{\partial x^4} - 2\frac{\partial g}{\partial x}\frac{\partial^3 f}{\partial x^3} - g\frac{\partial^4 f}{\partial x^4} \right)
\end{align*}

よって、
\begin{align*}
    (\text{与式}) &= \frac{\partial^3}{\partial x^3} \left( f\frac{\partial g}{\partial x} - g\frac{\partial f}{\partial x} \right) - 4r \frac{\partial}{\partial x} \left( f\frac{\partial g}{\partial x} - g\frac{\partial f}{\partial x} \right) - 2\frac{\partial r}{\partial x} \left( f\frac{\partial g}{\partial x} - g\frac{\partial f}{\partial x} \right) \\
    &= L \left( f\frac{\partial g}{\partial x} - g\frac{\partial f}{\partial x} \right)
\end{align*}
\end{proof}

\begin{story}
\begin{proof}[命題3.14の証明]
(3.82)に有理関数解 $w_l$ が存在するとして
\begin{equation*}
    L(w_l) + 2\frac{\partial r}{\partial t_l}(x,t) = 0
\end{equation*}
から，積分可能条件として
\[
\frac{\partial^{2} r}{\partial t_l\partial t_{h}}(x,t)=\frac{\partial^{2} r}{\partial t_h\partial t_{l}}(x,t)
\]
より、(wの交換関係は当たり前に仮定していて、そこからどういう数式が生まれるか見ている。)
\begin{align}
    0 &= \frac{\partial}{\partial t_h}L(w_l) - \frac{\partial}{\partial t_l}L(w_h)\\
    &= \left(\frac{\partial^{4}}{\partial t_h\partial x^{3}}-4\frac{\partial r}{\partial t_{h}}\frac{\partial}{\partial x}-4r\frac{\partial^{2}}{\partial t_{h}\partial x}-2\frac{\partial^{2}r}{\partial t_{h}\partial x}-2\frac{\partial r}{\partial t_{h}}\frac{\partial}{\partial x}\right)w_{l}\\
    &-\left(\frac{\partial^{4}}{\partial t_l\partial x^{3}}-4\frac{\partial r}{\partial t_{l}}\frac{\partial}{\partial x}-4r\frac{\partial^{2}}{\partial t_{l}\partial x}-2\frac{\partial^{2}r}{\partial t_{l}\partial x}-2\frac{\partial r}{\partial t_{l}}\frac{\partial}{\partial x}\right)w_{h}\\
    &= \left(\frac{\partial^{3}}{\partial x^{3}}-4r\frac{\partial}{\partial x}-2\frac{\partial r}{\partial x}\right)(\frac{w_{l}}{t_{h}}-\frac{w_{h}}{t_{l}})\\
    &= L\left( \frac{\partial w_l}{\partial t_h} - \frac{\partial w_h}{\partial t_l} \right) + L_{hl} \label{eq:3.84}
\end{align}
が成り立っている．ここで
\begin{equation*}
    L_{hl} = -4\frac{\partial r}{\partial t_h}\frac{\partial w_l}{\partial x} + 4\frac{\partial r}{\partial t_l}\frac{\partial w_h}{\partial x} - 2w_l\frac{\partial^2 r}{\partial x \partial t_h} + 2w_h\frac{\partial^2 r}{\partial x \partial t_l}
\end{equation*}
であるが，計算して補題3.6を使うと
\begin{align*}
    L_{hl} &= 2\frac{\partial w_l}{\partial x}L(w_h) - 2\frac{\partial w_h}{\partial x}L(w_l) + w_l\frac{\partial}{\partial x}L(w_h) - w_h\frac{\partial}{\partial x}L(w_l) \\
    &= L\left( w_l\frac{\partial w_h}{\partial x} - w_h\frac{\partial w_l}{\partial x} \right)
\end{align*}
したがって(3.84)より
\begin{equation*}
    L\left( \frac{\partial w_l}{\partial t_h} - \frac{\partial w_h}{\partial t_l} + w_l\frac{\partial w_h}{\partial x} - w_h\frac{\partial w_l}{\partial x} \right) = 0
\end{equation*}
となる．命題3.8の証明で見たように，仮定(P2), (P3)によれば，3階線型常微分方程式 $L(w)=0$ は非自明な有理関数解をもたない．すなわち
\begin{equation*}
    \frac{\partial w_l}{\partial t_h} - \frac{\partial w_h}{\partial t_l} + w_l\frac{\partial w_h}{\partial x} - w_h\frac{\partial w_l}{\partial x} = 0
\end{equation*}
以上の議論は，(3.83)が(3.82)の積分可能条件に他ならないことを示している．
\end{proof}
\end{story}
\end{document}